\chapter{The Days of the Week}
\label{ch:days}

Just as with the numbers, German goes its \textbf{Germanic} way here --- and
lands right next to English. The Romans named the weekdays for planet-gods
(\emph{Mars, Mercury, Jupiter, Venus}); the Germanic peoples swapped in
\emph{their own} gods, and German and English inherited that same set.
\emph{Practice lessons: \texttt{lessons/GE-C07-*.md}.}

\section{\emph{Montag} to \emph{Freitag} --- the gods of the Germanic week}
\label{lesson:wochentage-1}

\begin{sounds}
\emph{-tag} (``day'') is easy; watch \emph{Freitag} $=$ \emph{FRY-tahk}, with
the \emph{ei} $=$ English ``eye.''\\[3pt]
\emph{Mittwoch} ends in the raspy \emph{ch}: \emph{MIT-vokh}.
\end{sounds}

\begin{morphologybox}
\textbf{The pattern: [Germanic god] $+$ \emph{Tag} (``day'').} \textbf{Tag} $=$
``day'' (cousin of English \textbf{day}). Before it stands a god --- the
Germanic counterpart of the Roman planet.
\end{morphologybox}

\begin{center}
\begin{tabular}{@{}llll@{}}
\toprule
German & god & $=$ English & (Roman planet it replaced) \\
\midrule
\emph{Montag} & the \textbf{Moon} (\emph{Mond}) & \textbf{Mon}day & Moon \\
\emph{Dienstag} & \textbf{Tiw/Ziu}, the war-god & \textbf{Tues}day & Mars \\
\emph{Mittwoch} & \emph{(none --- ``mid-week'')} & \textbf{Wednes}day (Woden) & Mercury \\
\emph{Donnerstag} & \textbf{Donar} $=$ \textbf{Thor} (\emph{Donner} ``thunder'') & \textbf{Thurs}day & Jupiter \\
\emph{Freitag} & \textbf{Frija/Frigg}, love-goddess & \textbf{Fri}day & Venus \\
\bottomrule
\end{tabular}
\end{center}

\begin{cousinweb}
\textbf{\emph{Donnerstag} $\leftrightarrow$ Thursday.} \emph{Donner} is German
for ``thunder,'' and \emph{Donar} was the Germanic thunder-god --- the one the
Norse called \textbf{Thor}. So \emph{Donnerstag} $=$ ``Thor's day'' $=$
\textbf{Thursday}, standing in for the Roman \emph{Jupiter} (the sky-thunder
king). Same god-slot, three names.
\end{cousinweb}

\begin{culture}
\textbf{\emph{Mittwoch} is the odd one out.} It ought to be ``Wodan's day'' (as
English keeps \textbf{Wednes}day $=$ Woden's day). But the medieval Church
disliked naming a day for the chief pagan god, so German replaced it with plain
\textbf{Mittwoch}, ``\textbf{mid-week}'' --- a religious edit, just like
Portuguese numbering its days.
\end{culture}

\section{\emph{Samstag} and \emph{Sonntag} --- the German weekend}
\label{lesson:wochentage-2}

Two weekend days, and one more surprise: German's Saturday is \emph{not}
Saturn's day --- it is the \textbf{Sabbath}, sharing its root with the Romance
languages, not with English \emph{Saturday}. Sunday, though, stays the Sun's.

\begin{center}
\begin{tabular}{@{}llll@{}}
\toprule
German & $\leftarrow$ & meaning & note \\
\midrule
\emph{Samstag} & \emph{Sabbat} $\leftarrow$ Greek \emph{s\'{a}bbaton} $\leftarrow$ Hebrew \emph{shabb\={a}t} & the ``\textbf{Sabbath}'' & \emph{not} Saturn! \\
\emph{Sonntag} & \emph{Sun} $+$ \emph{Tag} (\emph{Sonne} ``sun'') & ``the \textbf{Sun's} day'' & $=$ English \textbf{Sun}day \\
\bottomrule
\end{tabular}
\end{center}

\begin{etymology}
\textbf{Samstag} surprises everyone: where English kept the Roman planet
(\textbf{Satur}day $=$ Saturn), German took the \textbf{Sabbath} --- Greek
\emph{s\'{a}bbaton} (from Hebrew \emph{shabb\={a}t}, ``to rest'') travelled up
through the early Church into German as \emph{Sabbat} $\to$ \emph{Samstag}. It
reached German through Greek-speaking Christians. So German \emph{Samstag} and
Spanish \emph{s\'{a}bado} are the \emph{same} Sabbath, while English alone
still points at Saturn.
\end{etymology}

\begin{culture}
In northern Germany you also hear \textbf{Sonnabend}, ``Sun-eve,'' the evening
before Sunday.\\[3pt]
\textbf{Sonntag} is simply ``the \textbf{Sun's} day,'' \emph{Sonne} $+$
\emph{Tag} --- the exact twin of English \textbf{Sun}day. Here Germanic and
English agree, and the Church did not rename it (unlike Romance, which made it
``the Lord's day,'' \emph{domingo}/\emph{dimanche}). It is a rare Germanic day
the Church left alone.
\end{culture}

All seven, in order: \emph{Montag, Dienstag, Mittwoch, Donnerstag, Freitag,
Samstag, Sonntag}. German names its days for \emph{gods}, like English, not for
planets --- with two edits: \emph{Mittwoch} (``mid-week,'' for Woden) and
\emph{Samstag} (the Sabbath, for Saturn). Next chapter builds on this toward
telling the time.
